\chapter{Controller Design and Implementation}
\label{ch:implementation}

This chapter describes the implementation of the Cartesian impedance controller derived in Chapter~2 for the Franka Emika Panda.
The implementation involves accessing model quantities through libfranka, evaluating the controller within the \(1\,\mathrm{ms}\) Franka Control Interface cycle, mapping the Cartesian wrench to joint torques, and logging the signals used later for evaluation.

\section{Real-Time Controller Structure}
\label{sec:implementation_scope}

The controller is implemented in C++ using libfranka and Eigen.
The libfranka library provides access to the robot state, the model interface, and the real-time torque-control callback.
Eigen is used for matrix and vector operations, such as computing Cartesian velocities, evaluating the impedance law, and mapping Cartesian quantities to joint torques.

The implementation is organized around the real-time callback function
\texttt{robot.control}.
This structure is required because the controller must compute and return a new
torque command within each \(1\,\mathrm{ms}\) control cycle.
Each experiment uses the same basic control-loop structure.
Only the desired end-effector pose and the stiffness and damping matrices are
changed between experiments.

The implementation follows four main requirements:

\begin{itemize}
    \item The measured robot state, including \(q\), \(\dot{q}\), and \(T_{\mathrm{EE}}\), is read at every control cycle.
    \item The Jacobian and Coriolis terms are obtained from the libfranka model interface.
    \item The Cartesian impedance wrench is computed using the selected stiffness and damping matrices.
    \item The resulting wrench is mapped to joint torques using \(J^\top(q)\), combined with the constrained-task null-space torque and the optional Coriolis term, and limited before being sent to the robot.
\end{itemize}

\subsection{Control-Loop Data Flow}
\label{sec:control_loop_structure}

Figure~\ref{fig:software_data_flow} shows the data flow of the implemented controller.
At each control step, the Franka Control Interface provides the current robot state, while the libfranka model interface provides the required kinematic and dynamic quantities.
The controller then computes the Cartesian error, evaluates the impedance law, maps the resulting wrench to joint torques, adds the projected null-space torque and the optional Coriolis term, and returns the final torque command.

\begin{figure}[h]
    \centering
    \begin{tikzpicture}[
        block/.style={draw, rounded corners, minimum width=4.2cm, minimum height=1.1cm, align=center},
        arrow/.style={->, thick}
    ]
        \node[block] (state) {Robot State\\ \(q,\dot{q},T_{\mathrm{EE}}\)};
        \node[block, below=0.9cm of state] (model) {Model Interface\\ \(J(q),\tau_c(q,\dot{q})\)};
        \node[block, below=0.9cm of model] (controller) {Impedance and Null-Space Controller\\ \(F=\begin{bmatrix} f\\ m \end{bmatrix},\ \tau_{\mathrm{null}}\)};
        \node[block, below=0.9cm of controller] (torque) {Torque Command\\
        \(\begin{aligned}
        \tau_{\mathrm{cmd}}={}&J^\top(q)\mathbf{F}+\tau_{\mathrm{null}}\\
        &+\tau_c(q,\dot{q})
        \end{aligned}\)};

        \draw[arrow] (state) -- (model);
        \draw[arrow] (model) -- (controller);
        \draw[arrow] (controller) -- (torque);
    \end{tikzpicture}
    \caption{Real-time data flow of the impedance controller implemented with libfranka.}
    \label{fig:software_data_flow}
\end{figure}

This structure follows the control law derived in Chapter~2.
This implementation chapter therefore focuses on obtaining the required
quantities from libfranka, evaluating the equations in real time, and ensuring
that the resulting torque command satisfies the safety limits.

\section{Cartesian Impedance Control Law}
\label{sec:state_model}

At the beginning of each control callback, the current joint configuration \(q\) and joint velocity vector \(\dot{q}\) are read from \texttt{franka::RobotState} and mapped to Eigen vectors.
The end-effector transformation is obtained from \texttt{state.O\_T\_EE}, which contains the homogeneous transformation matrix \(T_{\mathrm{EE}}\) of the libfranka end-effector/tool frame with respect to the robot base frame.
Thus, the position vector \(p_{\mathrm{EE}}\) used by the controller is the Cartesian position of this tool-frame origin, not an independently reconstructed DH point.

Although the theoretical forward kinematics can be constructed from the DH parameters, the implementation uses the kinematic quantities provided directly by libfranka.
The measured end-effector transformation is read from \texttt{state.O\_T\_EE}, and the Jacobian is obtained from the libfranka model interface.

The geometric Jacobian is obtained using

\begin{lstlisting}[language=C++]
model.zeroJacobian(franka::Frame::kEndEffector, state)
\end{lstlisting}

and is used for two purposes.
First, it maps joint velocities to Cartesian velocities,

\begin{equation}
\dot{x}
=
J(q)\dot{q}.
\label{eq:software_cartesian_velocity}
\end{equation}

Second, its transpose is used to map the Cartesian wrench to joint torques,

\begin{equation}
\tau_{\mathrm{task}}
=
J^\top(q)\mathbf{F}.
\label{eq:software_task_torque}
\end{equation}

The Coriolis term is also obtained from the libfranka model interface.
Together with the projected null-space torque, it is added after the task torque has been computed, resulting in the commanded torque

\begin{equation}
\tau_{\mathrm{cmd}}
=
J^\top(q)\mathbf{F}
+
\tau_{\mathrm{null}}
+
\tau_c(q,\dot{q}).
\label{eq:software_commanded_torque}
\end{equation}

An explicit gravity term is not added in this software equation.
This follows the structure of the official libfranka Cartesian impedance
example, where the user callback forms the commanded torque from the Cartesian
task torque, the null-space torque, and the Coriolis term, but does not add a
separate gravity vector.
The reason is that, for external joint-torque control, the Franka Control
Interface applies the desired torque command with robot-side gravity and
motor-friction compensation already included \cite{FrankaFCI2023}.
Therefore, the gravity compensation is not missing from the physical robot
control loop; it is handled below the user-level torque equation.
Adding another gravity vector in the user controller would compensate the same
effect twice and would not match the implementation structure used here.

This expression is the implementation form of the torque structure derived in
Chapter~2.
The purpose of this chapter is therefore to describe how the required
quantities are accessed, evaluated, and combined inside the real-time callback.

\section{Axis Constraint Mode}
\label{sec:axis_constraint_mode}

In addition to the standard Cartesian impedance behavior, the implemented
controller provides an axis constraint mode.
This mode defines which translational and rotational Cartesian directions are
actively regulated and which directions remain free for manual interaction.

The Cartesian position of the end-effector is

\begin{equation}
p
=
\begin{bmatrix}
x & y & z
\end{bmatrix}^{\top}.
\label{eq:axis_constraint_position}
\end{equation}

The position error is

\begin{equation}
\mathbf{e}_p
=
\mathbf{p}_d
-
\mathbf{p}_{\mathrm{EE}},
\label{eq:axis_constraint_position_error}
\end{equation}

where \(\mathbf{p}_d\) is the desired position and \(\mathbf{p}_{\mathrm{EE}}\) is the measured
end-effector position.
In the axis constraint mode, each translational axis can be fixed or free.
This is controlled by

\begin{equation}
\mathrm{fix}_p
=
\begin{bmatrix}
\mathrm{fix}_{p,x} &
\mathrm{fix}_{p,y} &
\mathrm{fix}_{p,z}
\end{bmatrix}.
\label{eq:fix_p_vector}
\end{equation}

If an axis is fixed, the desired position component is set to the initial
position.
If an axis is free, the desired position component is set to the current
measured end-effector position.
For each translational component \(i \in \{x,y,z\}\), the desired position is

\begin{equation}
p_{d,i}
=
\begin{cases}
p_{\mathrm{start},i}, & \text{if } \mathrm{fix}_{p,i}=1,\\
p_{\mathrm{EE},i}, & \text{if } \mathrm{fix}_{p,i}=0.
\end{cases}
\label{eq:axis_constraint_desired_position}
\end{equation}

As a result, free axes automatically have zero position error:

\begin{equation}
e_{p,i}
=
p_{d,i}
-
p_{\mathrm{EE},i}
=
0
\qquad
\text{if } \mathrm{fix}_{p,i}=0.
\label{eq:free_axis_position_error}
\end{equation}

The same idea is applied to the rotational error.
The orientation error is represented by

\begin{equation}
\mathbf{e}_R
=
\begin{bmatrix}
e_{R,x} &
e_{R,y} &
e_{R,z}
\end{bmatrix}^{\top}.
\label{eq:axis_constraint_orientation_error}
\end{equation}

The rotational constraint parameters are

\begin{equation}
\mathrm{fix}_R
=
\begin{bmatrix}
\mathrm{fix}_{R,x} &
\mathrm{fix}_{R,y} &
\mathrm{fix}_{R,z}
\end{bmatrix}.
\label{eq:fix_R_vector}
\end{equation}

If a rotational component is free, the corresponding component of the
orientation error is set to zero:

\begin{equation}
e_{R,i}
\leftarrow
\begin{cases}
e_{R,i}, & \text{if } \mathrm{fix}_{R,i}=1,\\
0, & \text{if } \mathrm{fix}_{R,i}=0.
\end{cases}
\label{eq:axis_constraint_orientation_mask}
\end{equation}

This allows the controller to regulate only the selected Cartesian directions
while leaving the remaining directions compliant or free.

The implementation follows this logic directly in the real-time callback.
The desired position is first initialized from the start pose.
Then, in axis constraint mode, every free translational component is replaced
by the current measured position.
For the rotational part, the selected free components of the rotation-vector
error are set to zero:

\begin{lstlisting}[style=cppstyle, language=C++, caption={Implementation of translational and rotational axis constraints.}, label={lst:axis_constraint_callback}]
Eigen::Vector3d p_d = p_start;
Eigen::Vector3d pdot_d = Eigen::Vector3d::Zero();

if (params.axis_constraint_mode) {
  // Fixed axis: desired position is the start value.
  // Free axis: desired position follows the measured position.
  p_d(0) = params.fix_p_x ? p_start(0) : p_EE(0);
  p_d(1) = params.fix_p_y ? p_start(1) : p_EE(1);
  p_d(2) = params.fix_p_z ? p_start(2) : p_EE(2);

  // Desired velocity is zero.
  // Free axes therefore have damping but no position spring.
  pdot_d(0) = 0.0;
  pdot_d(1) = 0.0;
  pdot_d(2) = 0.0;
}

Eigen::Vector3d e_p = p_d - p_EE;
Eigen::Vector3d e_R = orientationError(R_EE, R_d);

if (params.axis_constraint_mode) {
  // Fixed rotational components keep their orientation spring.
  // Free components have no rotational spring.
  e_R(0) = params.fix_R_x ? e_R(0) : 0.0;
  e_R(1) = params.fix_R_y ? e_R(1) : 0.0;
  e_R(2) = params.fix_R_z ? e_R(2) : 0.0;
}
\end{lstlisting}

Listing~\ref{lst:axis_constraint_callback} shows that the virtual plane is not
implemented as a separate geometric wall object.
It is created by the selected translational constraints.
For example, setting
\(\mathrm{fix}_p=[1\;0\;0]\) keeps the \(x\)-position close to the start
position while allowing motion in the \(y\)- and \(z\)-directions.
A tool orientation parallel to this virtual plane is obtained by choosing the
desired orientation \(R_d\) accordingly and keeping the required rotational
components active with \(\mathrm{fix}_R\).
This means that the controller keeps or regulates the orientation that is
defined in the experiment.
It does not automatically discover an unknown physical surface from contact and
then align the tool to that surface.
For a non-planar or arbitrary surface, the desired orientation would have to be
provided by the task definition or by an additional perception/contact
estimation method, which is outside the implemented controller.

\section{Virtual Plane and Virtual Line Behavior}
\label{sec:virtual_plane_line_behavior}

The axis constraint mode can be used to create virtual planes or virtual
lines.
A virtual plane is obtained by fixing one translational direction while
leaving the other two translational directions free.
For example,

\begin{equation}
\mathrm{fix}_p
=
\begin{bmatrix}
1 & 0 & 0
\end{bmatrix}
\label{eq:virtual_yz_plane_fixp}
\end{equation}

fixes the \(x\)-position while allowing motion in the \(y\)- and
\(z\)-directions.
This corresponds to a virtual \(y\)-\(z\) plane.
In this case, the desired position is

\begin{equation}
p_d
=
\begin{bmatrix}
x_{\mathrm{start}}\\
y_{\mathrm{EE}}\\
z_{\mathrm{EE}}
\end{bmatrix}.
\label{eq:virtual_yz_plane_desired_position}
\end{equation}

The resulting position error becomes

\begin{equation}
\mathbf{e}_p
=
\begin{bmatrix}
x_{\mathrm{start}} - x_{\mathrm{EE}}\\
0\\
0
\end{bmatrix}.
\label{eq:virtual_yz_plane_error}
\end{equation}

Thus, only deviations in the \(x\)-direction are corrected by the Cartesian
impedance controller.
The \(y\)- and \(z\)-directions remain free.

Similarly, fixing the \(y\)-axis creates a virtual \(x\)-\(z\) plane,

\begin{equation}
\mathrm{fix}_p
=
\begin{bmatrix}
0 & 1 & 0
\end{bmatrix},
\label{eq:virtual_xz_plane_fixp}
\end{equation}

and fixing the \(z\)-axis creates a virtual \(x\)-\(y\) plane,

\begin{equation}
\mathrm{fix}_p
=
\begin{bmatrix}
0 & 0 & 1
\end{bmatrix}.
\label{eq:virtual_xy_plane_fixp}
\end{equation}

A virtual line is generated by fixing two translational axes.
For example,

\begin{equation}
\mathrm{fix}_p
=
\begin{bmatrix}
1 & 1 & 0
\end{bmatrix}
\label{eq:virtual_z_line_fixp}
\end{equation}

fixes the \(x\)- and \(y\)-directions while allowing motion only in the
\(z\)-direction.
If all three translational axes are fixed,

\begin{equation}
\mathrm{fix}_p
=
\begin{bmatrix}
1 & 1 & 1
\end{bmatrix},
\label{eq:standard_cartesian_fixp}
\end{equation}

the controller behaves like a standard Cartesian impedance controller that
returns the end-effector to the initial position.

\subsection{Tool Orientation for a Given Surface}
\label{sec:tool_orientation_given_surface}

If the surface is known in advance, the tool orientation can be defined from
the surface geometry and then used as the desired orientation \(R_d\) of the
orientation controller.
This is different from automatically detecting an unknown surface during
contact.
In this case, the surface information is assumed to be given by the task
definition, for example as a surface normal vector.

Let

\begin{equation}
n_s
\in
\mathbb{R}^{3},
\qquad
\|n_s\|=1,
\label{eq:given_surface_normal}
\end{equation}

be the unit normal vector of the given surface, expressed in the base frame.
For a planar surface, \(n_s\) is constant.
For a curved surface, \(n_s\) can be interpreted as the local normal at the
current contact or reference point.
To make the tool frame parallel to the surface, one axis of the end-effector
frame is aligned with this normal.
In this work, the tool approach direction is represented by the
\(z_{\mathrm{EE}}\)-axis.
Therefore, a natural choice is

\begin{equation}
z_d
=
n_s .
\label{eq:desired_tool_z_axis}
\end{equation}

The remaining two axes of the desired orientation must lie in the tangent
plane of the surface.
Let \(a\in\mathbb{R}^{3}\) be an auxiliary direction that is not parallel to
\(n_s\).
A first tangent direction can be computed by projecting \(a\) onto the tangent
plane and normalizing it:

\begin{equation}
x_d
=
\frac{
a
-
n_s
\left(
n_s^{\top}a
\right)}
{\left\|
a
-
n_s
\left(
n_s^{\top}a
\right)
\right\|}.
\label{eq:surface_tangent_x_axis}
\end{equation}

The second tangent direction is then obtained from the cross product

\begin{equation}
y_d
=
z_d
\times
x_d .
\label{eq:surface_tangent_y_axis}
\end{equation}

The desired orientation matrix is formed from these three orthonormal axes:

\begin{equation}
R_d
=
\begin{bmatrix}
x_d & y_d & z_d
\end{bmatrix}.
\label{eq:desired_orientation_from_surface}
\end{equation}

With this construction, the desired tool orientation is no longer chosen
arbitrarily.
It is calculated from the given surface normal.
The controller can then use the same orientation error computation as in the
standard implementation:

\begin{equation}
\mathbf{e}_R
=
\mathrm{orientationError}
\left(
R_{\mathrm{EE}},R_d
\right).
\label{eq:surface_orientation_error}
\end{equation}

If all rotational axes are regulated, the controller tries to keep the complete
tool orientation equal to \(R_d\).
If only selected rotational components are regulated through
\(\mathrm{fix}_R\), then only those components of the orientation error remain
active.
This makes it possible to combine a given surface orientation with the
existing axis constraint mode.
For example, the translational constraint can define the virtual plane, while
the desired orientation \(R_d\) can keep the tool frame parallel to the given
surface.

For the axis-aligned virtual planes used in this implementation, the surface
normal is especially simple.
A virtual \(y\)-\(z\) plane has normal

\begin{equation}
n_s
=
\begin{bmatrix}
1 & 0 & 0
\end{bmatrix}^{\top},
\label{eq:yz_plane_surface_normal}
\end{equation}

while a virtual \(x\)-\(z\) plane has normal
\([0\;1\;0]^{\top}\), and a virtual \(x\)-\(y\) plane has normal
\([0\;0\;1]^{\top}\).
For a general non-axis-aligned surface, the same method can still be used if
the surface normal is provided by the task description, a CAD model, a
trajectory planner, or an external perception/contact-estimation method.
The present controller does not estimate this normal by itself; it only uses
the desired orientation that is supplied to it.

\subsection{Virtual-Center Self-Alignment Extension}
\label{sec:self_alignment_virtual_center}

The orientation construction in Section~\ref{sec:tool_orientation_given_surface}
defines how the tool can be kept parallel to a given surface.
However, the task description also requires self-aligning rotational
compliance with a virtual center of rotation that can be selected on, above,
or below the surface.
This requires an additional geometric coupling between translation and
rotation.
The purpose of the virtual center is to define the point about which the tool
appears to rotate while it aligns with the surface.

Let \(p_s\in\mathbb{R}^{3}\) be a point on the given surface and let
\(n_s\in\mathbb{R}^{3}\) be the unit surface normal.
The virtual center of rotation is defined as

\begin{equation}
p_c
=
p_s
+
h n_s,
\label{eq:virtual_center_from_surface}
\end{equation}

where \(h\in\mathbb{R}\) is the chosen offset from the surface.
The three important cases are

\begin{equation}
\begin{array}{lll}
h=0, & p_c \text{ on the surface},\\[2mm]
h>0, & p_c \text{ above the surface},\\[2mm]
h<0, & p_c \text{ below the surface}.
\end{array}
\label{eq:virtual_center_offset_cases}
\end{equation}

The vector from the virtual center to the end-effector is

\begin{equation}
r_c
=
p_{\mathrm{EE}}
-
p_c .
\label{eq:implementation_virtual_center_lever_arm}
\end{equation}

If the end-effector rotates by a small angle
\(\delta\theta\in\mathbb{R}^{3}\) about this virtual center, the corresponding
linear displacement of the end-effector is approximated by

\begin{equation}
\delta p_{\mathrm{rot}}
\approx
\delta\theta
\times
r_c .
\label{eq:implementation_vcr_position_motion}
\end{equation}

This equation is the key difference between simple orientation regulation and
self-alignment with a virtual center of rotation.
With simple orientation regulation, the controller tries to reduce the
orientation error at the end-effector frame.
With a virtual center of rotation, a rotation also implies a translational
motion determined by the lever arm \(r_c\).
The tool therefore behaves as if it pivots about the selected point \(p_c\).

The simplified coupled error used for such a controller can be written as

\begin{equation}
\mathbf{e}_c
=
\begin{bmatrix}
\mathbf{e}_p
-
\mathbf{e}_R
\times
\mathbf{r}_c\\
\mathbf{e}_R
\end{bmatrix},
\label{eq:implementation_vcr_coupled_error}
\end{equation}

where \(\mathbf{e}_p=\mathbf{p}_d-\mathbf{p}_{\mathrm{EE}}\) is the translational error and \(\mathbf{e}_R\) is
the orientation error computed with respect to the surface-based desired
orientation \(R_d\).
The upper part of \(\mathbf{e}_c\) contains the translational error corrected by the
motion that would be produced by rotating around the virtual center.
The lower part contains the rotational error itself.

The translational impedance force is then computed using the coupled
translational error

\begin{equation}
\mathbf{f}_c
=
K_p
\left(
\mathbf{e}_p
-
\mathbf{e}_R
\times
\mathbf{r}_c
\right)
+
D_p
\left(
\dot{\mathbf{p}}_d
-
\dot{\mathbf{p}}_{\mathrm{EE}}
\right).
\label{eq:implementation_vcr_force}
\end{equation}

The rotational impedance moment is computed as in the standard controller,

\begin{equation}
\mathbf{m}_R
=
K_R \mathbf{e}_R
-
D_R\omega_{\mathrm{EE}} .
\label{eq:implementation_vcr_rotational_moment}
\end{equation}

However, because the translational force acts at the end-effector while the
virtual center is located at \(\mathbf{p}_c\), this force also creates an additional
moment through the lever arm \(r_c\):

\begin{equation}
m_{\mathrm{lever}}
=
r_c
\times
f_c .
\label{eq:implementation_vcr_lever_moment}
\end{equation}

The wrench used by the controller is therefore

\begin{equation}
F_c
=
\begin{bmatrix}
f_c\\
m_R + r_c\times f_c
\end{bmatrix}.
\label{eq:implementation_vcr_corrected_wrench}
\end{equation}

This is the corrected virtual-center coupling used for the practical
Cartesian impedance structure.
It is simpler than a full spatial impedance formulation
\cite{Fasse1997SpatialImpedance}, but it includes the
two important effects: the orientation error changes the effective
translational error, and the translational force generates a moment about the
selected virtual center.
The resulting wrench is mapped to joint torques with the Jacobian transpose in
the same way as the standard Cartesian impedance controller.

For the self-aligning surface task, the desired orientation \(R_d\) is
constructed from the given surface normal, as described in
Section~\ref{sec:tool_orientation_given_surface}.
The rotational stiffness can be chosen lower around the axes that should align
compliantly, while the translational stiffness regulates the constrained
surface direction.
During contact, the tool can then rotate toward the surface orientation while
the selected virtual center determines whether this rotation occurs as if the
pivot point were on, above, or below the surface.

Thus, the specified self-aligning surface behavior consists of two parts.
First, the surface normal defines the desired tool orientation.
Second, the virtual center \(p_c\) defines the point about which the compliant
rotation is generated.
Only combining both parts gives the intended self-aligning behavior.
The axis-constraint mode alone can define a virtual plane and regulate a given
orientation, but a full self-aligning surface controller requires the
additional virtual-center coupling described in this section.

In relation to the implemented controller, the missing part is therefore not an
automatic discovery of an unknown surface.
For the task formulation considered here, the surface is assumed to be a given
task geometry, described by a point \(p_s\) and a normal vector \(n_s\).
The missing part is the complete online use of the selectable virtual center
\(p_c=p_s+h n_s\) inside the real-time control law.
To fully realize the specified behavior, the implementation would need to add
the parameter \(h\), compute \(p_c\) and \(r_c\) at every control step, replace
the independent position/orientation error by the coupled error
\eqref{eq:implementation_vcr_coupled_error}, and use the corrected
virtual-center wrench
\eqref{eq:implementation_vcr_corrected_wrench} in the torque command.
If the surface were not given in advance, an additional estimation layer would
be necessary to determine \(p_s\) and \(n_s\) from contact forces, external
torque estimates, vision, or another sensor source.
That surface-estimation layer is a different problem and is not part of the
implemented controller.

\section{Constrained-Task Jacobian for Null-Space Control}
\label{sec:active_jacobian_construction}

The impedance torque is always computed with the full geometric Jacobian
\(J(q)\in\mathbb{R}^{6\times7}\). The row-reduced Jacobian introduced here is
used only for the null-space projector. This distinction is important: the
controller does not replace \(J(q)\) in the Cartesian impedance law.

When the axis-constraint mode releases some translational or rotational
directions, the null-space term should preserve only the directions that are
actually regulated. For this purpose the implementation forms a constrained
task Jacobian

\begin{equation}
J_c
=
S_cJ(q)
\in
\mathbb{R}^{m\times7},
\label{eq:implementation_constrained_jacobian}
\end{equation}

where \(S_c\) selects the \(m\) rows belonging to fixed Cartesian axes. For
example, if \(x\), \(R_x\), and \(R_y\) are fixed, then

\begin{equation}
J_c
=
\begin{bmatrix}
J_x\\
J_{R_x}\\
J_{R_y}
\end{bmatrix}
\in
\mathbb{R}^{3\times7}.
\label{eq:implementation_constrained_jacobian_example}
\end{equation}

The code variable \texttt{J\_active} denotes this same matrix \(J_c\). The
name is kept in the listing because it is the variable used in the program,
but the thesis interpretation is simply: \(J_c\) is the Jacobian of the
currently constrained Cartesian task.

\section{Implemented Null-Space Term}
\label{sec:nullspace_optimization_implementation}

The theoretical null-space and singular-value terms were derived in
Section~\ref{sec:nullspace_svd}. In the implementation, the projector is built
from \(J_c\), not from the full \(6\times7\) Jacobian:

\begin{equation}
N
=
I
-
J_c^\top
\left(
J_cJ_c^\top+\lambda^2I
\right)^{-1}
J_c.
\label{eq:implementation_constrained_projector}
\end{equation}

The null-space torque has two parts,

\begin{equation}
\tau_{\mathrm{null}}
=
-D_{\mathrm{null}}N\dot{q}
+
N\left(K_{\sigma}s\alpha n\right).
\label{eq:implementation_nullspace_torque}
\end{equation}

The first term damps redundant joint motion. The second term is the optional
smallest-singular-value improvement. The implementation uses the scalar
parameter \texttt{nullspace\_k\_sigma}, corresponding to
\(K_{\sigma}=k_{\sigma}I_7\). The return-to-start term was removed because it
conflicted with the desired free motion in the plane.

The final torque command is still the Cartesian impedance task torque plus
the projected null-space torque and the optional Coriolis term,

\begin{equation}
\tau_{\mathrm{cmd}}
=
J^\top(q)F_{\mathrm{cmd}}
+
\tau_{\mathrm{null}}
+
\tau_c.
\label{eq:implementation_final_torque_with_nullspace}
\end{equation}

The null-space torque is rate- and norm-limited before being sent to the
robot, so it remains a bounded secondary contribution.

\subsection{Code-Level Constrained Null-Space Optimization}
\label{sec:code_active_nullspace}

The following code excerpts show the implementation of the constrained-task
null-space optimization.
The null space is computed from the constrained-task Jacobian introduced in
Section~\ref{sec:active_jacobian_construction}; this matrix contains only the
rows corresponding to the constrained Cartesian axes.
For \(J_c\in\mathbb{R}^{m\times7}\), the damped projector is
formed from an \(m\times m\) matrix, but the resulting projector \(N\) remains
\(7\times7\). Therefore, the projected null-space torque is still a
seven-dimensional joint torque vector.

The singular-value part uses a full SVD of \(J_c\). In the code this matrix is
named \texttt{J\_active}. Since the Panda has seven joints, Eigen's
\texttt{matrixV().col(6)} selects the seventh right singular vector. This
vector is used as one local joint-space search direction \(n\). The controller
tests \(q+\alpha n\) and \(q-\alpha n\), recomputes the constrained-task
Jacobian for both candidates, and keeps the sign that gives the larger
smallest singular value. The resulting secondary torque is projected with
\(N\), norm-limited, and then added to the Cartesian task torque. It does not
replace the impedance task; it only damps and biases redundant joint motion.

\begin{lstlisting}[style=cppstyle, language=C++, caption={Construction of the constrained task Jacobian.}, label={lst:active_task_jacobian}]
Eigen::MatrixXd buildActiveTaskJacobian(
    const Eigen::Matrix<double, 6, 7>& J,
    const Parameters& params) {

  int rows = 0;
  if (params.fix_p_x) ++rows;
  if (params.fix_p_y) ++rows;
  if (params.fix_p_z) ++rows;
  if (params.fix_R_x) ++rows;
  if (params.fix_R_y) ++rows;
  if (params.fix_R_z) ++rows;

  Eigen::MatrixXd J_active(rows, 7);
  int r = 0;

  if (params.fix_p_x) {
    J_active.row(r++) = J.row(0);
  }
  if (params.fix_p_y) {
    J_active.row(r++) = J.row(1);
  }
  if (params.fix_p_z) {
    J_active.row(r++) = J.row(2);
  }
  if (params.fix_R_x) {
    J_active.row(r++) = J.row(3);
  }
  if (params.fix_R_y) {
    J_active.row(r++) = J.row(4);
  }
  if (params.fix_R_z) {
    J_active.row(r++) = J.row(5);
  }

  return J_active;
}
\end{lstlisting}

Listing~\ref{lst:active_task_jacobian} selects the rows of the full
\(6 \times 7\) Jacobian that correspond to the constrained translational and
rotational axes.
Rows 0, 1, and 2 correspond to translational \(x\), \(y\), and \(z\) motion,
whereas rows 3, 4, and 5 correspond to rotational motion about these axes.

The constrained-task Jacobian is then used to compute a damped null-space projector.
This projector maps joint-space quantities into the null space of the
constrained task.

\begin{lstlisting}[style=cppstyle, language=C++, caption={Damped null-space projector of the constrained task.}, label={lst:active_task_projector}]
Eigen::Matrix<double, 7, 7> activeTaskNullspaceProjector(
    const Eigen::MatrixXd& J_active,
    double damping_lambda) {

  Eigen::Matrix<double, 7, 7> I =
      Eigen::Matrix<double, 7, 7>::Identity();

  if (J_active.rows() == 0) {
    return I;
  }

  Eigen::MatrixXd A =
      J_active * J_active.transpose()
      + damping_lambda * damping_lambda
            * Eigen::MatrixXd::Identity(
                  J_active.rows(),
                  J_active.rows());

  Eigen::MatrixXd J_pseudo =
      J_active.transpose() * A.inverse();

  Eigen::Matrix<double, 7, 7> N =
      I - J_pseudo * J_active;

  return N;
}
\end{lstlisting}

The damping term \(\lambda^2 I\) improves numerical robustness near singular
or poorly conditioned configurations.
If no Cartesian axis is fixed, the whole joint space is treated as free and
the identity matrix is returned.

Since the null-space torque is a seven-dimensional joint torque vector, a
separate norm limiter is used for this term:

\begin{lstlisting}[style=cppstyle, language=C++, caption={Norm limitation of the seven-dimensional null-space torque.}, label={lst:nullspace_torque_norm_limit}]
Eigen::Matrix<double, 7, 1> limitJointTorqueVectorNorm(
    const Eigen::Matrix<double, 7, 1>& v,
    double max_norm) {

  if (max_norm <= 0.0) {
    return v;
  }

  const double norm = v.norm();

  if (norm > max_norm && norm > 1e-12) {
    return max_norm * v / norm;
  }

  return v;
}
\end{lstlisting}

The limiter in Listing~\ref{lst:nullspace_torque_norm_limit} scales the
complete vector only when its norm exceeds the allowed maximum.
The direction of the null-space torque is therefore preserved.

The main null-space torque calculation is shown in
Listing~\ref{lst:active_task_nullspace_sigma}.
The previously used return-to-start term is not included, because it conflicted
with the desired free motion in the virtual plane.

\begin{lstlisting}[style=cppstyle, language=C++, caption={Constrained-task null-space damping and sigma-min optimization.}, label={lst:active_task_nullspace_sigma}]
Eigen::Matrix<double, 7, 1> tau_nullspace;
tau_nullspace.setZero();

if (params.use_nullspace_optimization) {

  Eigen::Map<const Eigen::Matrix<double, 7, 1>>
      q_current(state.q.data());

  Eigen::MatrixXd J_active =
      buildActiveTaskJacobian(J, params);

  Eigen::Matrix<double, 7, 7> N =
      activeTaskNullspaceProjector(
          J_active,
          params.active_nullspace_lambda);

  tau_nullspace =
      -params.nullspace_damping * (N * dq);

  if (params.nullspace_k_sigma > 0.0) {

    Eigen::JacobiSVD<Eigen::MatrixXd> svd_active(
        J_active,
        Eigen::ComputeFullU | Eigen::ComputeFullV);

    Eigen::Matrix<double, 7, 1> n =
        svd_active.matrixV().col(6);

    if (n.norm() > 1e-9) {
      n.normalize();

      Eigen::Matrix<double, 7, 1> q_plus =
          q_current + params.nullspace_alpha * n;

      Eigen::Matrix<double, 7, 1> q_minus =
          q_current - params.nullspace_alpha * n;

      std::array<double, 7> q_plus_array =
          eigenQToArray(q_plus);

      std::array<double, 7> q_minus_array =
          eigenQToArray(q_minus);

      std::array<double, 42> J_plus_array =
          model.zeroJacobian(
              franka::Frame::kEndEffector,
              q_plus_array,
              state.F_T_EE,
              state.EE_T_K);

      std::array<double, 42> J_minus_array =
          model.zeroJacobian(
              franka::Frame::kEndEffector,
              q_minus_array,
              state.F_T_EE,
              state.EE_T_K);

      Eigen::Map<const Eigen::Matrix<double, 6, 7>>
          J_plus_full(J_plus_array.data());

      Eigen::Map<const Eigen::Matrix<double, 6, 7>>
          J_minus_full(J_minus_array.data());

      Eigen::MatrixXd J_plus_active =
          buildActiveTaskJacobian(J_plus_full, params);

      Eigen::MatrixXd J_minus_active =
          buildActiveTaskJacobian(J_minus_full, params);

      const double sigma_min_plus =
          J_plus_active.rows() > 0
              ? Eigen::JacobiSVD<Eigen::MatrixXd>(
                    J_plus_active,
                    Eigen::ComputeThinU | Eigen::ComputeThinV)
                    .singularValues()
                    .minCoeff()
              : 0.0;

      const double sigma_min_minus =
          J_minus_active.rows() > 0
              ? Eigen::JacobiSVD<Eigen::MatrixXd>(
                    J_minus_active,
                    Eigen::ComputeThinU | Eigen::ComputeThinV)
                    .singularValues()
                    .minCoeff()
              : 0.0;

      const double sigma_direction =
          (sigma_min_plus > sigma_min_minus) ? 1.0 : -1.0;

      tau_nullspace +=
          N * (
              params.nullspace_k_sigma
              * sigma_direction
              * params.nullspace_alpha
              * n
          );
    }
  }

  tau_nullspace =
      limitJointTorqueVectorNorm(
          tau_nullspace,
          params.nullspace_tau_max);
}
\end{lstlisting}

The sigma-min term compares two small candidate motions in positive and
negative null-space direction.
The direction that produces the larger smallest singular value of the
constrained-task Jacobian is selected.
The resulting contribution is projected with \(N\), so it affects the
constrained Cartesian task as little as possible.

Finally, the null-space torque is added to the Cartesian impedance torque in
the real-time callback.
The Coriolis term is obtained from the libfranka model interface and is added
depending on the selected parameter:

\begin{lstlisting}[style=cppstyle, language=C++, caption={Combination of Cartesian task torque, null-space torque, and Coriolis torque.}, label={lst:task_nullspace_coriolis_torque}]
Eigen::Matrix<double, 7, 1> tau_raw =
    tau_task + tau_nullspace;

if (params.use_coriolis) {
  tau_raw += coriolis;
}

Eigen::Matrix<double, 7, 1> tau_limited =
    limitTorqueRate(
        tau_raw,
        tau_previous,
        params.delta_tau_max);

tau_previous = tau_limited;
\end{lstlisting}

The torque-rate limitation is applied after the task, null-space, and optional
Coriolis terms have been combined.
This keeps the torque command smooth before it is returned to the robot.

\section{Torque Saturation and Safety Handling}
\label{sec:torque_saturation_safety}

Several safety mechanisms are implemented to keep the commanded robot motion
smooth and to reduce the risk of abrupt behavior.
First, the Cartesian force command is limited by a maximum force magnitude:

\begin{equation}
\left\|
f_{\mathrm{cmd}}
\right\|
\leq
f_{\max}.
\label{eq:force_command_limit}
\end{equation}

Similarly, the Cartesian moment command is limited by

\begin{equation}
\left\|
m_{\mathrm{cmd}}
\right\|
\leq
m_{\max}.
\label{eq:moment_command_limit}
\end{equation}

The null-space torque is also limited:

\begin{equation}
\left\|
\tau_{\mathrm{null}}
\right\|
\leq
\tau_{\mathrm{null,max}}.
\label{eq:nullspace_torque_limit}
\end{equation}

This limit is necessary because the null-space torque is a seven-dimensional
joint torque vector,

\begin{equation}
\tau_{\mathrm{null}}
\in
\mathbb{R}^{7}.
\label{eq:implementation_nullspace_torque_dimension}
\end{equation}

In addition to magnitude limits, a torque-rate limit is applied.
The change of each joint torque between two consecutive control cycles is

\begin{equation}
\Delta \tau_i
=
\tau_{i,k}
-
\tau_{i,k-1},
\label{eq:torque_rate_difference}
\end{equation}

with

\begin{equation}
\left|
\Delta \tau_i
\right|
\leq
\Delta \tau_{\max}.
\label{eq:torque_rate_limit}
\end{equation}

This prevents sudden torque jumps and improves the smoothness of the
controller.
Furthermore, the controller includes a startup confirmation step.
After the user presses Enter once, the robot first attempts automatic error
recovery if necessary, then applies the custom collision behavior thresholds,
and finally moves to the initial joint configuration.

During the control loop, the user can stop the experiment safely by entering
\texttt{e} followed by Enter.
The controller then exits the torque-control callback using a controlled
\texttt{MotionFinished} command.
Custom collision thresholds are applied through the Franka collision behavior
interface.
These thresholds define the contact and collision sensitivity of the robot
during the experiment.
They do not replace the internal safety mechanisms of the robot, but they allow
the experiment to be configured for controlled manual interaction.

\section{Data Logging}
\label{sec:data_logging}

For the evaluation of the controller, relevant data are recorded during the
experiment.
The logged quantities include the desired and measured end-effector position,
the position error, the orientation error, the commanded Cartesian wrench, and
the commanded joint torques.

Because the controller runs in a real-time torque-control loop with a sampling
time of approximately \(1\,\mathrm{ms}\), file operations should be kept outside the
time-critical callback whenever possible.
Writing directly to a file during the callback can introduce delays and may
disturb the timing of the control loop.

The preferred logging approach stores the measured and computed values in a
memory buffer during the control loop.
After the experiment has finished, the buffered data are written to a CSV file
outside the real-time callback.
The general logging procedure is:

\begin{enumerate}
    \item During each control cycle, compute the relevant values.
    \item Store the values in a memory buffer.
    \item Finish the control loop either after the selected experiment duration
    or after the user stop command.
    \item Write the complete buffer to a CSV file after the robot control loop
    has ended.
\end{enumerate}

This approach keeps the real-time callback lightweight and avoids unnecessary
file I/O during torque control.
The resulting CSV file can then be used for post-processing, plotting, and
comparing different controller configurations.

\subsection{Complete Real-Time Callback}
\label{sec:control_impl}

The implemented control loop follows the sequence below:

\begin{enumerate}
    \item Read the current joint state \(q\), \(\dot{q}\), and end-effector transformation \(T_{\mathrm{EE}}\).
    \item Compute the geometric Jacobian \(J(q)\).
    \item Compute the Cartesian velocity using \(\dot{x}=J(q)\dot{q}\).
    \item Compute the position and orientation errors.
    \item Evaluate the Cartesian impedance wrench.
    \item Map the Cartesian wrench to joint torques using \(J^\top(q)\).
    \item Compute and add the projected null-space torque \(\tau_{\mathrm{null}}\).
    \item Add the optional Coriolis compensation.
    \item Limit the torque command before returning it to libfranka.
\end{enumerate}

The core structure of the implementation is shown in Listing~\ref{lst:impedance_callback}.
The listing keeps only the operations needed to show the control flow.
Task-specific files define the desired end-effector pose, the table-contact
configuration, and the candidate impedance parameters used for a given test.
They specify the positional stiffness values \(K_{1,p}\), \(K_{2,p}\), and
\(K_{3,p}\), the rotational stiffness values \(K_{1,R}\), \(K_{2,R}\), and
\(K_{3,R}\), and the corresponding damping values.

\begin{lstlisting}[style=cppstyle, language=C++, caption={Core structure of the real-time impedance-control callback with CSV logging.}, label={lst:impedance_callback}]
#include <algorithm>
#include <array>
#include <cmath>
#include <fstream>
#include <iomanip>
#include <Eigen/Dense>
#include <Eigen/Geometry>

// Define function eigenToArray: convert an Eigen torque vector to a std::array.
std::array<double, 7> eigenToArray(
    const Eigen::Matrix<double, 7, 1>& tau) {

  std::array<double, 7> tau_array{};

  for (int i = 0; i < 7; ++i) {
    tau_array[i] = tau(i);
  }

  return tau_array;
}

// Define function limitTorques: limit each commanded joint torque.
Eigen::Matrix<double, 7, 1> limitTorques(
    const Eigen::Matrix<double, 7, 1>& tau) {

  const double tau_max = 87.0;  // Example limit value.

  Eigen::Matrix<double, 7, 1> tau_limited;

  for (int i = 0; i < 7; ++i) {
    tau_limited(i) = std::max(-tau_max, std::min(tau_max, tau(i)));
  }

  return tau_limited;
}

// Load the robot model from libfranka.
// The model is used to compute the Jacobian and Coriolis terms.
franka::Model model = robot.loadModel();

// Desired end-effector position.
Eigen::Vector3d p_d;

// Desired orientation angles in radians.
double roll_d;
double pitch_d;
double yaw_d;

// Compute the desired orientation R_d using the ZYX convention:
// R_d = Rz(yaw_d) * Ry(pitch_d) * Rx(roll_d).
Eigen::Matrix3d R_d =
    Eigen::AngleAxisd(yaw_d, Eigen::Vector3d::UnitZ()).toRotationMatrix() *
    Eigen::AngleAxisd(pitch_d, Eigen::Vector3d::UnitY()).toRotationMatrix() *
    Eigen::AngleAxisd(roll_d, Eigen::Vector3d::UnitX()).toRotationMatrix();

// Positional stiffness values along the desired end-effector x-, y-, and z-axes.
double K1_p;
double K2_p;
double K3_p;

// Rotational stiffness values about the desired end-effector x-, y-, and z-axes.
double K1_R;
double K2_R;
double K3_R;
double D1_p;
double D2_p;
double D3_p;
double D1_R;
double D2_R;
double D3_R;

// Diagonal positional stiffness matrix in the desired end-effector frame.
Eigen::Matrix3d Kp_EE =
    Eigen::Vector3d(K1_p, K2_p, K3_p).asDiagonal();

// Diagonal rotational stiffness matrix in the desired end-effector frame.
Eigen::Matrix3d KR_EE =
    Eigen::Vector3d(K1_R, K2_R, K3_R).asDiagonal();

// Diagonal damping matrices in the desired end-effector frame.
Eigen::Matrix3d Dp_EE =
    Eigen::Vector3d(D1_p, D2_p, D3_p).asDiagonal();

Eigen::Matrix3d DR_EE =
    Eigen::Vector3d(D1_R, D2_R, D3_R).asDiagonal();

// Express the stiffness and damping matrices in the base frame.
Eigen::Matrix3d Kp_base = R_d * Kp_EE * R_d.transpose();
Eigen::Matrix3d Dp_base = R_d * Dp_EE * R_d.transpose();

Eigen::Matrix3d KR_base = R_d * KR_EE * R_d.transpose();
Eigen::Matrix3d DR_base = R_d * DR_EE * R_d.transpose();

// Open CSV file for logging experimental data.
std::ofstream log_file("impedance_experiment_log.csv");

// Write CSV header.
log_file << "time,"
         << "p_EE_x,p_EE_y,p_EE_z,"
         << "p_d_x,p_d_y,p_d_z,"
         << "e_p_x,e_p_y,e_p_z,"
         << "e_R_x,e_R_y,e_R_z,"
         << "pdot_x,pdot_y,pdot_z,"
         << "omega_x,omega_y,omega_z,"
         << "f_x,f_y,f_z,"
         << "m_x,m_y,m_z,"
         << "tau_1,tau_2,tau_3,tau_4,tau_5,tau_6,tau_7"
         << "\n";

// Initialize experiment time.
double time = 0.0;

// Real-time torque-control callback:
// libfranka provides the current robot state at each control cycle, including
// q, dq, and T_EE. The callback computes and returns the joint torque vector tau.
robot.control([&](const franka::RobotState& state,
                  franka::Duration period) -> franka::Torques {

  // Map the current joint positions q from the robot state to an Eigen vector.
  Eigen::Map<const Eigen::Matrix<double, 7, 1>>
      q(state.q.data());

  // Map the current joint velocities dq from the robot state to an Eigen vector.
  Eigen::Map<const Eigen::Matrix<double, 7, 1>>
      dq(state.dq.data());

  // Compute the geometric Jacobian at the end-effector using the libfranka model.
  std::array<double, 42> jacobian_array =
      model.zeroJacobian(franka::Frame::kEndEffector, state);

  // Map the Jacobian array to a 6x7 Eigen matrix.
  Eigen::Map<const Eigen::Matrix<double, 6, 7>>
      J(jacobian_array.data());

  // Compute the Cartesian end-effector velocity from the joint velocities.
  Eigen::Matrix<double, 6, 1> xdot = J * dq;

  // Extract the linear velocity and angular velocity of the end-effector.
  Eigen::Vector3d pdot = xdot.head<3>();
  Eigen::Vector3d omega = xdot.tail<3>();

  // Map the current homogeneous end-effector transformation T_EE to an Eigen matrix.
  Eigen::Map<const Eigen::Matrix<double, 4, 4>>
      T_EE(state.O_T_EE.data());

  // Extract the current end-effector position p_EE and orientation R_EE.
  Eigen::Vector3d p_EE = T_EE.block<3, 1>(0, 3);
  Eigen::Matrix3d R_EE = T_EE.block<3, 3>(0, 0);

  // Compute the axis-constrained desired position and position error.
  Eigen::Vector3d p_d_current = p_d;
  Eigen::Vector3d pdot_d = Eigen::Vector3d::Zero();

  if (params.axis_constraint_mode) {
    p_d_current(0) = params.fix_p_x ? p_start(0) : p_EE(0);
    p_d_current(1) = params.fix_p_y ? p_start(1) : p_EE(1);
    p_d_current(2) = params.fix_p_z ? p_start(2) : p_EE(2);
  }

  Eigen::Vector3d e_p = p_d_current - p_EE;

  // Compute the relative rotation from desired orientation to current orientation.
  Eigen::Matrix3d dR = R_d.transpose() * R_EE;

  // Compute the rotation angle from the trace of the relative rotation matrix.
  double cos_phi = (dR.trace() - 1.0) / 2.0;
  cos_phi = std::min(1.0, std::max(-1.0, cos_phi));
  double phi = std::acos(cos_phi);

  // Compute the rotational error vector using the axis-angle representation.
  Eigen::Vector3d e_R;
  if (phi < 1e-6) {
    e_R.setZero();
  } else {
    e_R =
        phi / (2.0 * std::sin(phi)) *
        Eigen::Vector3d(dR(2, 1) - dR(1, 2),
                        dR(0, 2) - dR(2, 0),
                        dR(1, 0) - dR(0, 1));
  }

  if (params.axis_constraint_mode) {
    e_R(0) = params.fix_R_x ? e_R(0) : 0.0;
    e_R(1) = params.fix_R_y ? e_R(1) : 0.0;
    e_R(2) = params.fix_R_z ? e_R(2) : 0.0;
  }

  // Compute the positional impedance force.
  Eigen::Vector3d f =
      Kp_base * e_p + Dp_base * (pdot_d - pdot);

  // Compute the rotational impedance moment.
  Eigen::Vector3d m =
      KR_base * e_R - DR_base * omega;

  // Combine the Cartesian force and moment into one 6D wrench.
  Eigen::Matrix<double, 6, 1> F;
  F.head<3>() = f;
  F.tail<3>() = m;

  // Map the Cartesian wrench to joint torques using the Jacobian transpose.
  Eigen::Matrix<double, 7, 1> tau_task =
      J.transpose() * F;

  // Compute the constrained-task null-space torque.
  // The detailed implementation is shown in the null-space listings above.
  Eigen::Matrix<double, 7, 1> tau_nullspace;
  tau_nullspace.setZero();

  // Read the Coriolis and centrifugal compensation torque tau_c from the libfranka model.
  std::array<double, 7> coriolis_array = model.coriolis(state);

  // Map the compensation term to an Eigen vector.
  Eigen::Map<const Eigen::Matrix<double, 7, 1>>
      tau_c(coriolis_array.data());

  // Compute the final commanded joint torque vector.
  Eigen::Matrix<double, 7, 1> tau =
      tau_task + tau_nullspace + tau_c;

  // Limit the commanded torques before sending them to the robot.
  Eigen::Matrix<double, 7, 1> tau_limited = limitTorques(tau);

  // Update elapsed experiment time.
  time += period.toSec();

  // Log measured and computed quantities for later evaluation.
  log_file << std::fixed << std::setprecision(6)
           << time << ","
           << p_EE(0) << "," << p_EE(1) << "," << p_EE(2) << ","
           << p_d_current(0) << "," << p_d_current(1) << "," << p_d_current(2) << ","
           << e_p(0) << "," << e_p(1) << "," << e_p(2) << ","
           << e_R(0) << "," << e_R(1) << "," << e_R(2) << ","
           << pdot(0) << "," << pdot(1) << "," << pdot(2) << ","
           << omega(0) << "," << omega(1) << "," << omega(2) << ","
           << f(0) << "," << f(1) << "," << f(2) << ","
           << m(0) << "," << m(1) << "," << m(2) << ","
           << tau_limited(0) << "," << tau_limited(1) << ","
           << tau_limited(2) << "," << tau_limited(3) << ","
           << tau_limited(4) << "," << tau_limited(5) << ","
           << tau_limited(6)
           << "\n";

  // Convert the limited Eigen torque vector to std::array and return it to libfranka.
  return franka::Torques(eigenToArray(tau_limited));
});
\end{lstlisting}

Listing~\ref{lst:impedance_callback} summarizes the main operations performed during one real-time control cycle.
The callback reads the current robot state, evaluates the Cartesian impedance law, adds the constrained-task null-space torque and the optional Coriolis term, limits the resulting torque command, and returns it to libfranka.
If a test configuration uses only translational compliance, the rotational
control part can be disabled by setting \(K_{R,0}=0_{3\times3}\) and
\(D_{R,0}=0_{3\times3}\). For the table-contact alignment experiment,
however, both translational and rotational components are relevant because the
contact wrench contains both force and moment terms.

The callback also writes the measured and computed quantities to a CSV file.
The logged values include the end-effector position, desired position, positional error, rotational error, Cartesian velocity, Cartesian force and moment, and limited joint torques.
This implementation structure ensures that all quantities required by the controller are evaluated within the real-time callback.
The torque command is computed from the current robot state at every control cycle and is limited before being returned to libfranka, which helps maintain safe operation during the experiments.

